% !TEX_{0} encoding = UTF-8 Unicode
\documentclass[11pt]{amsart}
\usepackage[top=30truemm,bottom=30truemm,left=30truemm,right=30truemm]{geometry} 
\usepackage{amsmath, amssymb, array}
\usepackage[all]{xy}
\usepackage{ascmac}
\usepackage[dvips]{graphicx}
\usepackage{color}
\usepackage{amscd}
\usepackage[driverfallback=dvipdfm]{hyperref}
%\usepackage[dvipdfmx]{hyperref}
\usepackage{amsrefs}
\usepackage{cleveref}
\usepackage{shuffle}
\usepackage{mathrsfs}
\usepackage{listings}
\usepackage{comment}

\newfont{\cyr}{wncyr10 scaled\magstep0}



\pagestyle{plain}



\newcommand{\Q}{\mathbb{Q}}
\newcommand{\Z}{\mathbb{Z}}

\DeclareMathOperator{\Tr}{Tr} %trace

\newcommand{\bb}[1]{\boldsymbol{#1}}
\newcommand{\ba}{\boldsymbol{a}}



%\text{\'et}



\theoremstyle{plain}
 \newtheorem{theorem}{Theorem}[section]
 \crefname{theorem}{Theorem}{Theorems}
 \newtheorem{proposition}[theorem]{Proposition}
 \crefname{proposition}{Proposition}{Propositions}
 \newtheorem{lemma}[theorem]{Lemma}
 \crefname{lemma}{Lemma}{Lemmas}
 \newtheorem{corollary}[theorem]{Corollary}
 \crefname{corollary}{Corollary}{Corollaries}
 \newtheorem{conjecture}[theorem]{Conjecture}
 \crefname{conjecture}{Conjecture}{Conjectures}
 \newtheorem{hypothesis}[theorem]{Hypothesis}
 \crefname{hypothesis}{Hypothesis}{Hypotheses}
 \newtheorem{question}[theorem]{Question}
 \crefname{question}{Question}{Questions}
 \newtheorem{problem}[theorem]{Problem}
 \crefname{problem}{Problem}{Problems}

 
\theoremstyle{definition} 
 \newtheorem{definition}[theorem]{Definition}
 \crefname{definition}{Definition}{Definitions}
 \newtheorem{example}[theorem]{Example}
 \crefname{example}{Example}{Examples}
 \newtheorem{remark}[theorem]{Remark}
 \crefname{remark}{Remark}{Remarks}





\title{A note on the Diophantine equation $2n^{2} = 1+q+ \cdots +q^{\alpha}$ and application to odd perfect numbers}
\author{Yoshinosuke Hirakawa}
\address[Yoshinosuke Hirakawa]{Department of Mathematics \\ Faculty of Science and Technology \\ Tokyo University of Science, 2641, Yamazaki, Noda, Chiba, Japan}
\email{hirakawa\_yoshinosuke@rs.tus.ac.jp \\ hirakawa@keio.jp}
	\thanks{This research was supported by KAKENHI 21K13779.}
	\subjclass[2010]{primary 11A41; %Primes
	secondary
		11A05; %Multiplicative structure; Euclidean algorithm; greatest common divisors
	%	11B83; %Special sequences and polynomials
		11D41; %Higher degree equations; Fermat's equation
	%	11D57; %Multiplicative and norm form equations
	%	11D61; %Exponential Diophantine equations
	%	11E76; %Forms of degree higher than two
	%	11N13; %Primes in congruence classes
	%	11N32; %Primes represented by polynomials; other multiplicative structure of polynomial values
		11R11; %Quadratic extensions
	%	11R16; %Cubic and quartic extensions
		11R27; %Units and factorization
		11G30; %Curves of arbitrary genus or genus ≠1 over global fields 
		}
\keywords{perfect numbers; multiply perfect numbers; Diophantine equations}
%\date{\today}
 


\begin{document}



%%%Title&Author

\maketitle



%%% Abstract 
\begin{abstract}
Let $N$ be an odd perfect number. Then, Euler proved that there exist some integers $n, \alpha$ and a prime $q$ such that $N = n^{2}q^{\alpha}$, $q \nmid n$, and $q \equiv \alpha \equiv 1 \bmod 4$.
In this note, we prove that the ratio $\sigma(n^{2})/q^{\alpha}$ cannot be a square unless $\alpha = 1$.
It is a direct consequence of a certain property of the Diophantine equation $2n^{2} = 1+q+ \cdots +q^{\alpha}$, whose proof is based on the prime ideal factorization in the quadratic order $\mathbb{Z}[\sqrt{1-q}]$. We give also a slight generalization to odd multiply perfect numbers.
\end{abstract}



%\tableofcontents



%%%Contents


%%%
\section{Main result}

Let $N$ be a perfect number,
i.e., $\sigma(N) = 2N$,
where $\sigma(N)$ is the sum of divisors of $N$.
Euler proved that
the even perfect numbers $N$ correspond bijectively to the Mersenne primes $q = 2^{p}-1$
as follows:
\[
	N
	= 2^{p-1}q.
	%= (2^{\frac{p-1}{2}})^{2}q
	%= \frac{q+1}{2} \cdot q^{1}
\]
Historically, the fact that such $N$ is a perfect number was already known in Euclide's era.

In what follows,
we assume that $N$ is an odd perfect number.
Then, Euler proved that
there exist some integers $n, \alpha$ and a prime $q$ such that
\[
	N = n^{2}q^{\alpha},
	\quad q \nmid n,
	\quad \text{and}
	\quad q \equiv \alpha \equiv 1 \bmod 4.
\]
For its proof, see e.g.\ \cite[Theorem 2.2]{Chen-Luo}.
Since the quantity
\[
	\gcd(n^{2}, \sigma(n^{2}))
	= \gcd\left( \frac{\sigma(q^{\alpha})}{2} \cdot \frac{\sigma(n^{2})}{q^{\alpha}}, \frac{\sigma(n^{2})}{q^{\alpha}} \right)
	= \frac{\sigma(n^{2})}{q^{\alpha}}
\]
is a counterpart of the quantity
\[
	\gcd(2^{p-1}, \sigma(2^{p-1}))
	= \frac{\sigma(2^{p-1})}{q}
	= 1
\]
for even $N$,
it is natural to ask whether $\sigma(n^{2}) = q^{\alpha}$ also for odd $N$ (e.g.\ \cite{Suryanarayana}).
The purpose of this note is to give an answer to this question in a more precise form as follows.



%%%
\begin{theorem} \label{2-perfect}
Suppose that $N = n^{2}q^{\alpha}$ is an odd perfect number as above.
\begin{enumerate}
\item
If $\alpha = 1$, then $\sigma(n^{2})/q > 1$.  

\item
If $\alpha > 1$, then $\sigma(n^{2})/q^{\alpha} $ is not a square.
\end{enumerate}
In any case, $\sigma(n^{2}) \neq q^{\alpha}$.
\end{theorem}

%Since we do not know whether $\gcd(\sigma(n^{2})/q^{k}, \sigma(q^{k})) = 1$ or not, $\sigma(n^{2})/q^{\alpha} \neq 1$ does not imply that it is not a square.



\begin{remark} \label{DHP1}
In \cite[Theorem 1]{Dandapat-Hunsucker-Pomerance},
Dandapat, Hunsucker, and Pomerance proved that
$\sigma(n^{2}) \neq q^{\alpha}$ for every $\alpha$
with some variants for multiply perfect numbers (see \cref{DHP}).
Thus, \cref{2-perfect} gives a refinement of \cite{Dandapat-Hunsucker-Pomerance}
at least partially.
\end{remark}



%%%
\section{Proof}

\subsection{The case of $\alpha = 1$}

\cref{2-perfect} for $\alpha = 1$ is well-known.
Let $\omega(n)$ be the number of distinct prime divisors of $n$
and $n = \prod_{i = 1}^{\omega(n)} p_{i}^{\alpha_{i}}$ with primes $p_{i}$ and $\alpha_{i} > 0$.
Then, the equality $q = \sigma(n^{2})$ implies that
\[
	q = \prod_{i = 1}^{\omega(n)} (1+p_{i}+\cdots+p_{i}^{2\alpha_{i}}),
\]
and hence $\omega(n) = 1$.
However, it contradicts the known lower bound e.g.\ $\omega(n) \geq 9$ due to Nielsen \cite[Theorem 3.8]{Nielsen}.

In fact,
we can deduce a huge lower bound for $\sigma(n^{2})/q$
from known results on the divisors of $n$;
for example, a single inequality $\omega(n) \geq 9$ implies that
\[
	\frac{\sigma(n^{2})}{q}
	\geq \sigma(3^{2}) \cdot \sigma(5^{2}) \cdot \sigma(7^{2}) \cdot \sigma(11^{2})
		\cdot \sigma(13^{2}) \cdot \sigma(17^{2}) \cdot \sigma(19^{2}) \cdot \sigma(23^{2})
	= 36163554870725919.
\]
For more discussion, see e.g.\ \cite{Nielsen}.



\subsection{The case of $\alpha > 1$}

%In order to prove the statement for $\alpha > 1$ by contradiction,
Suppose that $\sigma(n^{2})/q^{\alpha}$ is a square $n_{0}^{2}$.
Then, since $2(n/n_{0})^{2} = \sigma(q^{\alpha})$,
\cref{2-perfect} is a consequence of the following property
of the Diophantine equation $2n^{2} = \sigma(q^{\alpha})$.


\begin{theorem} \label{main}
There exist no triples $(\alpha, n, q)$ of integers such that
$\alpha > 1$, $q \equiv 1 \bmod 4$ is a prime, and $2n^{2} = \sigma(q^{\alpha})$.
\end{theorem}


\begin{proof}
We prove it by contradiction.
Suppose that $2n^{2} = \sigma(q^{\alpha})$.
We may assume that $\alpha \geq 3$ is odd.
In this case, the above equation can be rephrased as follows:
\[
	2(q-1)n^{2}
	%= (q-1) \cdot \frac{q^{k+1}-1}{q-1}
	= (q^{\frac{\alpha+1}{2}}-1)(q^{\frac{\alpha+1}{2}}+1).
\]
Since $q-1$ divides $q^{\frac{\alpha+1}{2}}-1$, $2$ divides $q^{\frac{\alpha+1}{2}}+1$, and 
\[
	\gcd(q^{\frac{\alpha+1}{2}}-1, q^{\frac{\alpha+1}{2}}+1)
	= \gcd(q^{\frac{\alpha+1}{2}}-1, 2)
	= 2,
\]
there exist some (coprime) divisors $n_{1}, n_{2}$ of $n$ such that 
\[
	\begin{cases}
	(q-1)n_{1}^{2} = q^{\frac{\alpha+1}{2}}-1, \\
	2n_{2}^{2} = q^{\frac{\alpha+1}{2}}+1.
	\end{cases}
\]
In particular, the former equation implies that
\[
	(1+n_{1}\sqrt{1-q})(1-n_{1}\sqrt{1-q})
	= q^{\frac{\alpha+1}{2}}
	= (1+\sqrt{1-q})^{\frac{\alpha+1}{2}}(1-\sqrt{1-q})^{\frac{\alpha+1}{2}}.
\]
Since $1 \pm n_{1}\sqrt{1-q}$ are coprime in the ring $\Z[\sqrt{1-q}]$
and $1 \pm \sqrt{1-q}$ are prime elements of $\Z[\sqrt{1-q}]$ above $q$,
the uniqueness of the prime ideal decomposition outside the ideals dividing $1-q$ (see e.g.\ \cite[\S7]{Cox})
implies that
there exists some non-negative integer $v$ such that
\[
	1+n_{1}\sqrt{1-q}
	= \epsilon(1+\sqrt{1-q})^{v}(1-\sqrt{1-q})^{\frac{\alpha+1}{2}-v},
\]
where $\epsilon \in \Z[\sqrt{1-q}]^{\times} = \{ \pm1 \}$.
Moreover,
if $0 < v < \frac{\alpha+1}{2}$,
then $1+n_{1}\sqrt{1-q}$ must be divisible by $q$,
a contradiction.
%which contradicts the fact that $1$ (and $n_{1}$) cannot be divisible by $q$.
Therefore, $v = 0$ or $\frac{\alpha+1}{2}$.
In any case,
we obtain the following equality
\[
	2
	= \epsilon\Tr_{\Q(\sqrt{1-q})/\Q}((1 \pm \sqrt{1-q})^{\frac{\alpha+1}{2}})
	%= \epsilon((1 + \sqrt{1-q})^{\frac{\alpha+1}{2}} + (1 - \sqrt{1-q})^{\frac{\alpha+1}{2}}) \\
	= 2\epsilon + 2\epsilon(1-q)\sum_{i = 1}^{\left[ \frac{\alpha+1}{4} \right]} \binom{\frac{\alpha+1}{2}}{2i} 1^{\frac{\alpha+1}{2}-2i} (1-q)^{i-1},
	%= 2\epsilon + 2\epsilon\binom{\frac{\alpha+1}{2}}{2}(1-q)^{i} + 2\epsilon\sum_{i = 2}^{\left[ \frac{\alpha+1}{4} \right]} \binom{\frac{\alpha+1}{2}}{2i} 1^{\frac{\alpha+1}{2}-2i} (1-q)^{i},
\]
which implies that
\[
	1 \equiv \epsilon \bmod q-1,
	\quad \text{i.e.,} \quad
	\epsilon = 1.
\]
Therefore,
by combining it with the assumption that $\alpha \geq 3$,
we obtain the following identity
\[
	0
	= 1 + \sum_{i = 2}^{\left[ \frac{\alpha+1}{4} \right]} \frac{\binom{\frac{\alpha+1}{2}}{2i}}{\binom{\frac{\alpha+1}{2}}{2}} (1-q)^{i-1}
	= 1 + \sum_{i = 2}^{\left[ \frac{\alpha+1}{4} \right]} \frac{\binom{\frac{\alpha-3}{2}}{2i-2}}{2i-1} \cdot \frac{(1-q)^{i-1}}{i}.
\]
However,
the right hand side is a $2$-adic unit whenever $q \equiv 1 \bmod 4$,
a contradiction.
\end{proof}



In the above proof,
we obtain the following by-product.

\begin{theorem} \label{by-product}
There exist no triples $(\beta, n, q)$ of integers such that
$\beta \geq 1$, $q \equiv 1 \bmod 4$ is a prime, and $n^{2} = \sigma(q^{\beta})$.
\end{theorem}



%%%
\section{Application to multiply perfect numbers}

A positive integer $N$ is called $k$-perfect if $\sigma(N) = kN$.
The classical perfect numbers are $2$-perfect in this sense.
In \cite{Chen-Luo}, Chen and Luo generalized Euler's theorem on odd $2$-perfect numbers to odd $k$-perfect numbers with even $k$ as follows.
In what follows, $v_{2}(k)$ denotes the integer $v$ such that $2^{v} \mid k$ but $2^{v+1} \nmid k$.

\begin{theorem}[{cf.\ \cite[Theorem 2.2]{Chen-Luo}} \label{Chen-Luo}
\footnote{
	Since the statement of \cite[Theorem 2.2]{Chen-Luo} is incorrect,
	the author recovered a correct statement from its proof.
	}
]
Let $k \geq 2$ be an even integer and $N$ be an odd $k$-perfect number.
Then, there exist an integer $s$ such that $1 \leq s \leq v_{2}(k)$,
some integers $n, \alpha_{1}, \dots, \alpha_{s}$,
and distinct primes $q_{1}, \dots, q_{s}$ such that
\[
	N = n^{2}q_{1}^{\alpha_{1}} \cdots q_{s}^{\alpha_{s}}
	\quad \text{and}
	\quad q_{i} \nmid n.
\]
Moreover,
there exist some non-negative integers $a_{1}, \dots, a_{s}, b_{1}, \dots, b_{s}$ such that
\[
	v_{2}(k) - s = \sum_{i = 1}^{s} a_{i} + \sum_{i = 1}^{s} b_{i},
	\quad p_{i} \equiv 2^{a_{i}+1}-1 \bmod 2^{a_{i}+2},
	\quad \text{and}
	\quad \alpha_{i} \equiv 2^{b_{i}+1}-1 \bmod 2^{b_{i}+2}.
\]
\end{theorem}

\begin{corollary} \label{Chen-Luo2}
Let $k$ be an integer such that $k \equiv 2 \bmod 4$ and $N$ be an odd $k$-perfect number.
Then, there exist some integers $n, \alpha$ and a prime $q$ such that
\[
	N = n^{2}q^{\alpha},
	\quad q \nmid n,
	\quad q \equiv \alpha \equiv 1 \bmod 4.
\]
\end{corollary}



Thanks to \cref{Chen-Luo2},
we can prove the following by the same argument as \cref{2-perfect}.

%%%
\begin{theorem} \label{k-perfect}
In the setting of \cref{Chen-Luo2},
further assume that $k/2$ is a square.
\begin{enumerate}
\item
If $\alpha = 1$, then $\sigma(n^{2})/q > 1$. %i.e., $\gcd(kn^{2}, \sigma(n^{2})) \geq 3$

\item
If $\alpha > 1$, then $\sigma(n^{2})/q^{\alpha}$ is not a square.
\end{enumerate}
In any case, $\sigma(n^{2}) \neq q^{\alpha}$.
\end{theorem}

\begin{proof}
For $\alpha = 1$, we use the lower bound $\omega(n) \geq k^{2}-2$ due to McCarthy \cite[Corollary]{McCarthy}.
For $\alpha > 1$, it is sufficient to note that
\[
	2 \cdot \frac{\frac{k}{2} \cdot n^{2}}{\frac{\sigma(n^{2})}{q^{\alpha}}} = \sigma(q^{\alpha}),
\]
which contradicts \cref{main} if $k/2$ and $\sigma(n^{2})/q^{\alpha}$ are squares.
\end{proof}



Here, it should be mentioned that
the result \cite[Theorem 1]{Dandapat-Hunsucker-Pomerance} of Dandapat, Hunsucker, and Pomerance can be restated as follows.

\begin{theorem} [{\cite[Corollary]{Dandapat-Hunsucker-Pomerance}}] \label{DHP}
Let $N = mq^{\alpha}$ be a $k$-perfect number with an integer $m$ and a prime $q$ such that$\sigma(m) = q^{\alpha}$ and  $q \nmid m$.
Then, $N = 672$ or an even perfect number.
\end{theorem}

Thus, by combining \cref{DHP} and \cref{Chen-Luo2},
we see that $\sigma(n^{2}) \neq q^{\alpha}$ for every odd $k$-perfect number $N$ with $k \equiv 2 \bmod 4$.
In this view,
\cref{k-perfect} is a refinement of this result under the condition that $N$ is odd and $k/2$ is a square.



\appendix
\section*{Acknowledgement}
The author would like to thank Takato Suzuki for valuable discussion on perfect numbers.


 
\begin{bibdiv}
\begin{biblist}
\bibselect{perfect}
\end{biblist}
\end{bibdiv}



\end{document}